\chapter{总结与展望}

针对Peach 模糊测试框架因缺乏协议语义感知能力而导致的变异盲目、难以触发深层逻辑的问题，本研究提出了一种协议语义感知变异增强方法，并将其集成于 Peach 框架中，开发了 LLMPeach 模糊测试工具。该工具使用由大模型智能体端到端生成的数据模型、变异器和修复器，采用兼顾变异深度与随机性的两阶段变异策略，力求在保持输入结构合理性的同时，更有针对性地对协议字段语义进行探索。

通过在多个协议实现上的实验评估，本研究证实了 LLMPeach 能有效提升被测协议实现的代码分支覆盖率，并能够触及传统 Peach 难以覆盖的深层逻辑路径。这表明结合大语言模型的语义感知机制生成的测试组件能够显著弥补传统模糊测试在语义维度的不足，提升复杂网络协议的漏洞挖掘效率。

尽管 LLMPeach 在协议模糊测试方面取得了初步的成功，但仍存在一些局限性和改进空间。
首先，当前的智能体设框架较为简单，未来可以探索更复杂的交互式智能体设计，以提升测试组件生成的质量，并降低开销。其次，当前的测试组件生成中只考虑了数据模型，未来可以进一步引入协议状态机的建模与生成，以实现更为全面的协议语义感知。另外，两阶段变异策略中的参数选择也值得进一步研究和优化。
最后，当前的实验评估主要集中在分支覆盖率等静态指标上，未来可以结合漏洞挖掘效率、实际漏洞发现数量等更具实用意义的指标进行评估，以更全面地验证方法的有效性。